\chapter*{\hfill{\centering  ВВЕДЕНИЕ}\hfill}

\addcontentsline{toc}{chapter}{ВВЕДЕНИЕ}

При работе со словами, необходимо их сравнивать, причем необходима конкретная метрика, которая покажет, насколько посимвольно одно слово отличается от другого. Это актуально, например, для корректировки опечаток в тектовых редакторах и приложениях, использующих поиск по строке.

Одной из таких метрик является Расстояние Левенштейна.
Данное расстояние является метрикой, измеряющей по модулю разность между двумя последовательностями символов \cite{levenshtein}.

\label{sec:targets}
Целью данной лабораторной работы является описание и исследование алгоритмов поиска расстояний Левенштейна и Дамерау---Левенштейна.

Для поставленной цели необходимо выполнить следующие задачи:
\begin{enumerate}
	\item Рассмотреть следующие алгоритмы поиска:
	\begin{enumerate}
		\item линейный алгоритм поиска расстояния Левенштейна;
		\item рекурсивный алгоритм поиска расстояния Левенштейна;
		\item рекурсивный алгоритм поиска расстояния Левенштейна c кешированием;
		\item линейный алгоритм поиска расстояния Дамерау---Левенштейна.
	\end{enumerate}
	\item Написать программу, реализующую вышеперечисленные алгоритмы поиска.
	\item Провести исследование затрачиваемого процессорного времени и памяти при различных реализациях алгоритмов.
	\item Провести сравнительный анализ алгоритмов.
\end{enumerate}